\section{결 론}
\label{sec:conclusion}

본 논문에서는 \code{HAETAE} 전체 서명 경로의 오류주입 지점을
상류 연산, 응답 연산 및 거부 판정으로 구분하고, SW 오류
에뮬레이션으로 각 의미론적 오류가 만드는 취약 조건을 분석하였다.
응답에서 비밀키 곱 항이 제거되는 T1과 일회용 난수 항이 제거되는
T2에 대해서는 $\svec_1$ 복원 관계를 도출하였다. T2는 단일 오류
응답과 공개 챌린지를 이용하며, T1은 동일 nonce의 정상 응답과
결함 응답 사이의 차분을 이용한다.

SW T2 실험에서는 30회 중 25회 목표 오류 응답이 방출되었고,
각 응답에서 $\svec_1$의 768개 계수가 모두 일치하였다. 본 논문의
핵심 물리 결과인 T1 실험에서는 곱 결과 버퍼를 사전 초기화한
인과대조 구현에 실제 클럭 글리치를 주입하였다. 공격 구간을 식별한
후 최대 900회 중 108회에서 정상 기준 응답과 두 결함 응답으로
$\svec_1$ 768/768을 복원하였고, 별도 실행에서도 594회 이내에
동일한 결과를 재현하였다.

한편 미수정 융합 응답 구현에 총 650회의 물리 클럭 글리치를 주입한
비교 실험에서는 깨끗한 T2가 관측되지 않았다. T1과 T2의 대비는
오류의 실제 효과가 컴파일된 load--compute--store 구조, 반복 범위와
버퍼 초기 상태에 의해 달라질 수 있음을 실험적으로 시사한다. 따라서
격자 기반 전자서명의 오류주입 안전성은 의미론적 오류가 만드는
대수적 영향과 특정 구현에서의 물리 실현성을 구분하여 평가해야 한다.

물리 T1 실험은 고정 nonce, 인과대조 빌드 및 내부 계측을 사용하므로
미수정 구현의 종단 공격으로 해석하지 않는다. 직렬화 서명만을 이용한
복원, 공개키 기반 후보 검증과 위조는 후속 공격 고도화로 남긴다.
또한 RB의 분포 누설 정량화와 다른 컴파일 구조 및 주입 수단에 대한
평가가 필요하다. 대응 측면에서는 응답 데이터 흐름과 서명 경계를
보호하는 경량 자체검증 방향을 후속 연구로 확장하고자 한다.
